\chapter{结果与分析}
\section{结果分析与局限}

当前系统已经实现了面向比赛提交和现场演示的 HarmonyOS 透明泡泡渲染效果。主要结果包括：

\begin{itemize}
  \item 通过多阶段 FBO 近似双界面折射，背景参照物和天空盒能在泡泡中产生明显扭曲。
  \item 两种薄膜虹彩模式已经在 Native 层实现，界面开放 Kim 与 Airy 两种展示模式，便于说明实时公式与多波长近似之间的差异。
  \item 动态膜厚让虹彩随时间缓慢流动，不依赖静态纹理。
  \item 触摸交互能够产生局部膜厚增强和波纹，配合参数面板形成更自然的现场展示反馈。
  \item 主泡泡接入 DBSTT/vortex sheet 风格表面动态，16 个多展示泡泡通过距离排序与 Alpha 混合合成，用于丰富空间构图与比赛演示层次。
  \item ArkTS 页面、NAPI 接口、Native 渲染层和可执行 HAP 程序已经形成完整的提交链路。
\end{itemize}

同时，系统也有明确局限：

\begin{itemize}
  \item 屏幕空间折射不能表现屏幕外物体，也不能严格模拟多泡泡之间的相互折射。
  \item Belcour Airy 模式是实时 shader 中的项目级简化，不包含完整微表面模型和解析光谱预积分。
  \item 动态膜厚主要由噪声、排液趋势和触摸扰动驱动，没有求解完整的薄膜厚度流体方程。
  \item DBSTT 模拟只用于单个闭合泡泡，不支持原论文中的泡沫拓扑变化和 Plateau 边界。
  \item 展示泡泡只做视觉漂移，不参与主泡泡的物理碰撞或流体耦合。
\end{itemize}

如果继续迭代比赛版本，后续优先考虑补强多泡泡之间的相互折射、碰撞与排序一致性，并进一步完善移动端性能记录。

\section{总结}

本项目从透明肥皂泡的主要视觉现象出发，实现了一个面向赛题二提交的 HarmonyOS 渲染系统。渲染方面，项目通过多阶段 FBO 合成近似双界面折射，并用低分辨率 FBO 与坐标修正兼顾移动端性能和多泡泡显示稳定性。光学方面，项目实现了 Kim2012和Belcour Airy 两种薄膜虹彩模式，能够展示实时公式计算、预计算纹理和多波长 Airy 近似之间的差异。动态方面，项目结合程序化膜厚变化、触摸波纹、DBSTT 风格主泡泡形变和多泡泡风吹漂移，使最终结果更符合比赛对“效果、交互、完整性”的要求。

后续如继续扩展该项目，可进一步加入更严格的光谱到 RGB 积分、更真实的膜厚流体模拟、多泡泡之间的相互折射和碰撞，以及基于体积或拓扑变化的泡沫模拟。
